\section{Results}
\label{sec:results}

We present results organized by research question, with each finding interpreted in terms of our claims.

\subsection{Main Results: SA-Correctness Correlation (RQ1)}

Our core finding: \textbf{pylint score achieves $r=0.87$ correlation with pass@1}, far exceeding the $r\geq0.35$ threshold for moderate predictive power.

\begin{table}[t]
\centering
\caption{SA metric correlations with pass@1. Bold indicates primary result. Mypy omitted due to numerical artifact.}
\label{tab:main-results}
\begin{tabular}{@{}lcccc@{}}
\toprule
SA Metric & $r_{\text{raw}}$ & $p$-value & $r_{\text{partial}}$ (LOC) & $p$-value \\
\midrule
\textbf{pylint\_score} & \textbf{0.868} & $3.8 \times 10^{-129}$ & \textbf{0.873} & $\mathbf{2.9 \times 10^{-132}}$ \\
radon\_cc & $-0.494$ & $2.4 \times 10^{-27}$ & $-0.569$ & $2.4 \times 10^{-37}$ \\
mypy\_errors & --- & --- & --- & --- \\
\bottomrule
\end{tabular}
\end{table}

\textbf{Key Observations:}

\begin{enumerate}
\item \textbf{Pylint achieves strong correlation ($r=0.87$, $p<10^{-132}$).} This far exceeds our 0.35 threshold. Higher pylint scores reliably predict test-passing code---the core claim is validated with high confidence.

\item \textbf{LOC control has minimal effect ($r_{\text{raw}}=0.868 \rightarrow r_{\text{partial}}=0.873$).} The correlation actually \emph{increases} slightly after controlling for code length, indicating the SA signal is genuine and not an artifact of shorter code being both cleaner and more correct.

\item \textbf{Radon cyclomatic complexity shows moderate negative correlation ($r=-0.57$).} As expected, lower complexity predicts correctness. This provides a second independent signal, though weaker than pylint.
\end{enumerate}

\subsection{Ensemble Analysis (RQ2)}

\textbf{Surprising finding: ensemble combination degrades performance.}

The weighted ensemble (pylint + radon) achieves $r=0.86$, \emph{below} pylint alone ($r=0.87$).

\begin{table}[h]
\centering
\caption{Ensemble vs.\ individual metric correlation.}
\label{tab:ensemble}
\begin{tabular}{@{}lcc@{}}
\toprule
Configuration & $r_{\text{ensemble}}$ & $p$-value \\
\midrule
Pylint only & 0.873 & $2.9 \times 10^{-132}$ \\
\textbf{Optimal ensemble (90\% pylint, 10\% radon)} & \textbf{0.861} & $1.1 \times 10^{-124}$ \\
50/50 ensemble & 0.664 & --- \\
\bottomrule
\end{tabular}
\end{table}

\textbf{Interpretation:} The ensemble degradation indicates pylint and radon capture \emph{redundant} rather than \emph{orthogonal} quality dimensions. Adding radon introduces noise without new signal. This is a ``less is more'' result: practitioners need only pylint for effective filtering.

\textbf{RQ2 Gate: FAIL} --- Ensemble does not outperform individual metrics. However, this negative result is practically useful: it simplifies deployment.

\subsection{Cross-Model Generalization (RQ3)}

\textbf{Finding: Correlation generalizes across all LLMs, but with higher variance than specified.}

\begin{table}[t]
\centering
\caption{Per-model pylint-correctness correlation. All models significant at $p<0.001$.}
\label{tab:cross-model}
\begin{tabular}{@{}lccc@{}}
\toprule
Model & $r_{\text{partial}}$ (pylint) & $p$-value & Significant \\
\midrule
\textbf{GPT-4} & \textbf{0.424} & $7.0 \times 10^{-8}$ & \checkmark \\
Claude-3 & 0.860 & $8.4 \times 10^{-45}$ & \checkmark \\
CodeLlama & 0.845 & $8.8 \times 10^{-42}$ & \checkmark \\
Codestral & 0.859 & $1.4 \times 10^{-44}$ & \checkmark \\
\midrule
\textbf{Mean} & \textbf{0.747} & --- & --- \\
\textbf{Std} & \textbf{0.186} & --- & --- \\
\bottomrule
\end{tabular}
\end{table}

\textbf{Key Observations:}

\begin{enumerate}
\item \textbf{All four models show statistically significant correlation ($p<0.001$).} The finding generalizes: SA-correctness relationship is not model-specific.

\item \textbf{Three models cluster tightly ($r=0.84$--$0.86$).} Claude-3, CodeLlama, and Codestral show nearly identical correlation strength.

\item \textbf{GPT-4 is an outlier ($r=0.42$).} Still significant and above threshold, but notably lower than other models. This inflates cross-model variance.

\item \textbf{Variance exceeds threshold ($\text{std}=0.19 > 0.15$).} The GPT-4 outlier causes gate failure on the strict variance criterion.
\end{enumerate}

\textbf{RQ3 Gate: FAIL (technical)} --- Variance exceeds 0.15. However, the directional finding holds: all models show $r>0.35$, confirming generalization in principle.

\subsection{Summary of Gate Outcomes}

\begin{table}[h]
\centering
\begin{tabular}{@{}lllll@{}}
\toprule
Hypothesis & Gate & Criterion & Result & Interpretation \\
\midrule
H-E1 & MUST\_WORK & $\geq$95\% valid rate & \textbf{PASS} (100\%) & Infrastructure validated \\
H-M1 & MUST\_WORK & $r\geq0.35$, $p<0.05$ & \textbf{PASS} ($r=0.87$) & Core claim validated \\
H-M2 & SHOULD\_WORK & $r_{\text{ensemble}} > r_{\text{individual}}$ & \textbf{FAIL} & Ensemble unnecessary \\
H-C1 & SHOULD\_WORK & $\text{std}(r) < 0.15$ & \textbf{FAIL} ($\text{std}=0.19$) & Generalizes with variance \\
\bottomrule
\end{tabular}
\end{table}
